% Path: docs/qualitative_section_template.tex
% Replace the placeholder filenames below with the actual figure files
% you want to discuss in the final report.

\section{Qualitative Analysis}

This section should compare representative outputs from the baseline model
and from the best-performing fine-tuned model, and should also include at
least one failure case.

\begin{figure}[H]
    \centering
    % Example:
    % \includegraphics[width=0.75\textwidth]{../figures/baseline_cpu/vis_X_filename.png}
    \fbox{\parbox[c][5cm][c]{0.75\textwidth}{\centering Insert baseline qualitative figure here}}
    \caption{Qualitative example from the baseline model. Replace this placeholder with a real figure from \texttt{figures/baseline\_cpu/}.}
    \label{fig:baseline_qualitative}
\end{figure}

\begin{figure}[H]
    \centering
    % Example:
    % \includegraphics[width=0.75\textwidth]{../figures/ablation_finetune_cpu/vis_X_filename.png}
    \fbox{\parbox[c][5cm][c]{0.75\textwidth}{\centering Insert fine-tuned qualitative figure here}}
    \caption{Qualitative example from the fine-tuned model. Replace this placeholder with a real figure from \texttt{figures/ablation\_finetune\_cpu/}.}
    \label{fig:finetune_qualitative}
\end{figure}

\begin{figure}[H]
    \centering
    % Use a figure that shows a clear failure case, such as a generic caption,
    % an incorrect object, or an incomplete sentence.
    \fbox{\parbox[c][5cm][c]{0.75\textwidth}{\centering Insert failure case here}}
    \caption{Failure case used to discuss model limitations. Replace this placeholder with a real figure selected from the generated outputs.}
    \label{fig:failure_case}
\end{figure}

\noindent Suggested discussion points:
\begin{itemize}
    \item whether the generated caption identifies the main object correctly;
    \item whether the action or scene context is captured;
    \item whether the caption is generic, incomplete, or incorrect;
    \item whether the fine-tuned model is more specific than the baseline.
\end{itemize}
